\chapter{Введение}

\section{Цель}  

Построить сеть состоящую из коммутатора и компьютеров.
Разделить сеть на сегменты с возможностью передачи трафика между подсетями.
\section{Задачи}

\begin{enumerate}
\item Подключить компьютеры V7C1 и V7C2 к портам коммутатора Gi1/0/1 и
Gi1/0/2.
\item Подключить компьютеры V7C3 и V7C4 к портам коммутатора Gi1/0/11 и
Gi1/0/12.
\item Подключить порт коммутатора Gi1/0/6 к интерфейсу маршрутизатора Gi1/0/1.
\item Создать на коммутаторе VLAN 10 и 20. Добавить порты Gi1/0/1 и Gi1/0/2 в
VLAN 10. Добавить порты Gi1/0/11 и Gi1/0/12 в VLAN 20. Порты коммутатора Gi1/0/1,
Gi1/0/2, Gi1/0/11, Gi1/0/12 перевести в режим доступа.
\item На коммутаторе, порт Gi1/0/6 перевести в магистральный режим и разрешить
передвижение трафика с 10 и 20 VLAN.
\item На маршрутизаторе создать Sub - интерфейсы Gi1/0/1.10 и Gi1/0/1.20.
\item Интерфейсу Gi1/0/1.10 назначить I7 - адрес 192.168.10.1/24. Интерфейсу
Gi1/0/1.20 назначить I7 - адрес 172.16.20.1/24.
\item Компьютерам V7C1 и V7C2, назначить адреса из сети 192.168.10.0/24.
Компьютерам V7C3 и V7C4, назначить адреса из сети 172.16.20.0/24.
\end{enumerate}

\begin{figure}[H] 
    \centering
    \includegraphics[width=0.75\textwidth]{top.png}
    \caption{Пример собранной сети}
\end{figure}

\endinput